\section{Corrections, source audit, and integration safeguards}
\label{audit:section}

\subsection{What was inspected and what was not established}

The source inspection covered the canonical manuscript's status and
relevant harmonic, finite-part, Stieltjes, Gamma, and Gaussian chapters;
the incoming archives were extracted and their relevant theorems,
normalizations, future questions, and verification notes compared.
The most direct dependencies are the ordered-Hurwitz,
triple-Stieltjes--Tornheim, coincident-Stieltjes, directional-zeta, and
resonance-coordinate reports. Two additional pinned formal-algebra
sources supply the finite separator and the depth-preserving Cayley
projector. This is a focused mathematical audit of the material used
here, not a line-by-line certification of the entire multi-volume
manuscript or every incoming claim.

No false theorem was identified in the selected incoming cubic and
coincident claims. In particular, the two previously supplied forms of
the diagonal cubic Tornheim relation are compatible after the zeta
functional equation is applied. Their common third-derivative coordinate
does not disappear merely because a reflected cubic difference is
evaluated in Section~\ref{ct:section}. The older corrections to the
rejected $S_8$ coefficient vector and related proof-status warnings are
already reflected in the inspected manuscript; they should not be
presented as newly discovered errors.

There is one definite external formula correction, proved in the next
subsection. The remaining items below are precise safeguards for
integrating the present results. They are not allegations that the
inspected source texts already make each prohibited inference.

% Optional audit subsection; all labels/citations have hp: prefix.
\subsection{A sign correction and its full logarithmic extension}
\label{hp:sec:tin-audit}

For this comparison use the unshifted raw-power harmonic zeta function
\[
 \mathcal H_p(s)=\sum_{n\ge1}\frac{H_n^p}{n^s},\qquad p\ge0,
 \quad \mathcal H_0(s)=\zeta(s),
\]
and define its regular Laurent coefficients by
\begin{equation}\label{hp:eq:tin-convention}
 \mathcal H_p(1+u)=\operatorname{PP}_{u=0}\mathcal H_p(1+u)
       +\sum_{j\ge0}\frac{(-1)^j}{j!}
                         \widetilde\gamma_{p,j}u^j.
\end{equation}
Thus $\widetilde\gamma_{0,j}=\gamma_j$, the ordinary Stieltjes
constant, and $\widetilde\gamma_{p,j}$ corresponds to Tin's
${}_p\widetilde\gamma_j$ in equation (14) of
\cite{hp:Tin2025}.

\begin{proposition}[Convergent identities for every logarithmic moment]
\label{hp:prop:tin-all-log}
For all integers $p,j\ge0$,
\begin{equation}\label{hp:eq:tin-all-log}
 \boxed{\quad
 \sum_{n\ge1}\frac{H_n^p(H_n-\log n-\gamma)\log^j n}{n}
   =\widetilde\gamma_{p+1,j}
        -\widetilde\gamma_{p,j+1}
        -\gamma\widetilde\gamma_{p,j}.\quad}
\end{equation}
The series is an ordinary absolutely convergent series.
\end{proposition}
\begin{proof}
Since $H_n-\log n-\gamma=1/(2n)+O(n^{-2})$, the Dirichlet
series with numerator $H_n^p(H_n-\log n-\gamma)$ and every fixed
logarithmic derivative converge normally in $\Re s>0$. In $\Re s>1$
its sum equals
\[
 \mathcal H_{p+1}(s)+\mathcal H_p'(s)-\gamma\mathcal H_p(s),
\]
so every principal part cancels at $s=1$. Comparing the coefficient
of $u^j$ in \eqref{hp:eq:tin-convention}, and then multiplying by
$(-1)^j j!$, proves \eqref{hp:eq:tin-all-log}.
\end{proof}

With the convention \eqref{hp:eq:tin-convention}, the first term on
the right of equation (22) of version 1 of \cite{hp:Tin2025},
printed page 5, must be
$-{}_p\widetilde\gamma_1$, in place of the printed
$+{}_p\widetilde\gamma_1$. This corrects that sign; no other term
of that equation is being audited here. The correction agrees with
the paper's own $p=1$ equation (11). Indeed,
\[
 \widetilde\gamma_{1,0}=\frac{\gamma^2+\zeta(2)}2,
 \qquad
 \widetilde\gamma_{2,0}=\frac{\gamma^3+5\zeta(3)}3
\]
give the two illustrative cases
\begin{align}
 \sum_{n\ge1}\frac{H_n-\log n-\gamma}{n}
   &=\frac{\zeta(2)-\gamma^2}{2}-\gamma_1,
       \label{hp:eq:tin-pzero}\\
 \sum_{n\ge1}\frac{H_n(H_n-\log n-\gamma)}{n}
   &=\frac53\zeta(3)-\frac\gamma2\zeta(2)
       -\frac{\gamma^3}{6}-\widetilde\gamma_{1,1}.
       \label{hp:eq:tin-pone}
\end{align}
The two displayed regular constants follow directly from
$\sum_{n\le N}H_n/n=(H_N^2+H_N^{(2)})/2$ and
$\sum_{n\le N}H_n^2/n=H_N^3/3+
\sum_{n\le N}H_n/n^2-H_N^{(3)}/3$, together with
$\sum_{n\ge1}H_n/n^2=2\zeta(3)$.


\subsection{Five safeguards with mathematical consequences}

\paragraph{Retain the full resonant amplitude.}
In \eqref{ct:R}, for odd $N=p+r$, the counterterm contains the entire
polynomial $P_p(u)/u$. Replacing it by its principal pole alone changes
the constant and higher coefficients. The difference is observable
already in
\[
 \FP\int_0^1\psi'(x)K(x)\,dx=2\zeta'(2),\qquad
 \FP\int_0^1\psi(x)K'(x)\,dx=-2\zeta'(2)-2\zeta(2).
\]
The harmonic correction in \eqref{ct:zero-index} is part of the identity,
not a removable convention. A source implementation should store the
completed germ, rather than store its pole residue and reconstruct a
finite part heuristically.

\paragraph{A total derivative can have a nonzero finite-part integral.}
With the fixed endpoint coordinates,
\[
 \FP\int_0^1 f'(x)\,dx=\CT_{x\to1}f(x)-\CT_{x\to0}f(x)
\]
whenever the endpoint Laurent-log expansions exist. Thus
$\FP\int_0^1(\psi K)'=-2\zeta(2)$, although the Gamma primitive
in \eqref{ct:loggamma-primitive} has zero boundary constant difference.
These two cases must be checked separately. A blanket rule that all
regularized total derivatives integrate to zero would be false.

\paragraph{Coordinate cancellation is not arithmetic minimality.}
The locus $b=c=d$ in Section~\ref{oq:section} is exactly the locus
where the coefficients of the named additional coordinates in the
displayed normal form vanish. The three-ray reconstruction is an
invertible coordinate calculation. Neither statement proves that
$\eta,\rho,\kappa,T$ are arithmetically independent over a specified
field of ordinary constants. Additional functional relations could
alter a later minimal presentation. The proposed status language is
``proved spanning normal form with an exact cancellation locus in that
form.''

\paragraph{A finite part and a spectral derivative use different data.}
At $s=-m$, multiplication by $1/\Gamma(s)$ kills the holomorphic
Mellin remainder in the constant Laurent coefficient of the raw-power
series. It generally does not kill its contribution to the coefficient
of $s+m$. Consequently, the polynomial formula for $C_{p,m}(a)$
cannot be differentiated with respect to the discrete label $m$ to
obtain a spectral derivative. Actual argument differentiation is valid
and gives \eqref{hp:eq:transport}, including its residue term. This
distinction is essential for any primitive or higher-Stieltjes
extension.

\paragraph{The complete Cayley result still has a fixed additional family.}
The ideal $\mathcal J$ in Section~\ref{gf:sec} is unrestricted, but
$\mathscr R_4$ is the enumerated additional family. The new theorem
rules out every construction within
$\mathcal J_7^-+\operatorname{span}_{\mathbb Q}\mathscr R_4$.
It does not rule out additional double-shuffle, distribution, or other
functional information. Nor does it prove that the actual period
residual is nonzero or that every possible proof must pass through
depth five. The integer pairing is a formal algebraic certificate,
not a lower bound for a period.

\subsection{Proposed canonical status updates}

The following replacements are suitable for the research-status record.
\begin{enumerate}
\item Mark the depth-four independent-slope request as answered by a
spanning formula with convergent coordinates and a formal cancellation
locus. Leave arithmetic minimality and a classification modulo all
functional relations open.
\item Extend the distinct-seam reflected closure to arbitrary repeated
seams by the completed common-seam transform and finite rational
partial fractions. Keep the generic unreflected cubic reduction open.
\item Add raw harmonic powers with continuous outer shift as a separate
family from elementary-symmetric harmonic numerators. Record the
all-even centered cancellation and the finite-part polynomial theorem.
\item Replace the proposed ``complete Cayley symmetry and retry the
same bounded rows'' step by the proved nonmembership theorem. Keep
$S_6$ and the revised $S_8$ explicitly conjectural.
\item Add a version-specific note for Tin's equation (22), with
\eqref{hp:eq:tin-all-log} as the proof of the corrected sign. Do not
silently change the Laurent convention to conceal the discrepancy.
\end{enumerate}
